\section{March 16}

\text{Recall:}
Let $M$ be a smooth (resp. holomorphic) manifold of dimension $n$. 
The following are equivalent:
\begin{enumerate}
  \item involutive subbundle $\calF \subset T_X$ (i.e. $[\calF, \calF] \subset \calF$) or rank $k$;
  \item foliation (resp. holomorphic foliation) of rank $r$ on $M$.
\end{enumerate}


\subsection{Foliations in algebraic geometry}

  Let $\kk$ be a field, $X/\kk$ a smooth variety.
  Recall the (sheaf of algebraic) sections of the tangent bundle $\calT_X$ is defined as the sheaf determined by 
  \[ U \mapsto \Der_\kk(\calO_X(U)), \quad \forall \text{ affine open } U \subset X. \]
  % \Yang{why above definition is not a (pre)sheaf?}
  Then $\calT_X$ is a locally free sheaf of rank $n$.
  
  \begin{definition}
    Assume $\characteristic \kk = 0$. 
    A \emph{foliation} on $X$ is an involutive subbundle $\calF \subset \calT_X$.
  \end{definition}

  \begin{remark}
    There is no natural notion of algebraic leaf, but there is a ``formal'' notion.
  \end{remark}

  Denote by $\hat{X}_P$ the formal completion of $X$ at $P$.
  We have 
  \[\hat{X}_P \cong \Spf \kk[[x_1, \dots, x_n]], \]
  where $x_1, \dots, x_n$ are local coordinates at $P$ since $X$ is smooth.
  

  \begin{definition}[formal leaf]
    Assume $P \in X(\kk)$ is a $\kk$-rational point.
    A \emph{formal leaf} of $\calF$ through $P$ is a closed smooth formal subscheme $\hat{L}_P \subset \hat{X}_P$ such that $\calF|_{\hat{L}_P} = \Image(\calT_{\hat{L}_P} \to \calT_{\hat{X}_P}|_{\hat{L}_P})$.
  \end{definition}

  \begin{theorem}
    Let $\calF$ be a foliation on $X$ and $P \in X(\kk)$ a $\kk$-rational point. 
    Then there exists a unique formal leaf $\hat{L}_P$ of $\calF$ through $P$.
  \end{theorem}

  \begin{remark}
    The differential equation have solution by formal power series.
    The construction can be found in \cite{Bost}.
  \end{remark}

  \begin{remark}
    If $\kk = \bbC$, then $X/\bbC$ gives a complex manifold $X^{\an}$, and $\calF$ gives a holomorphic foliation $\calF^{\an}$ on $X^{\an}$.
    Then the formal leaf $\hat{L}_P$ is the formal completion of the unique holomorphic leaf $L_P$ of $\calF^{\an}$ through $P$.
  \end{remark}

  \begin{question}[Fundamental Problem]
    When does there exist an (smooth) algebraic subvariety $L_P \subset X$ such that $\calF|_{L_P} = \Image(\calT_{L_P} \to \calT_X|_{L_P})$?

    Or equivalently, when does the formal leaf $\hat{L}_P$ algebraize? 
    (i.e. there exists an algebraic subvariety $L_P \subset X$ such that $\hat{L}_P$ is the formal completion of $L_P$ at $P$)
  \end{question}

  Now assume $\characteristic \kk = p > 0$.
  There is two structures on $\calT_X$:
  \begin{itemize}
    \item the Lie bracket $[-, -]$;
    \item the $p$-th power map $(-)^{[p]}: \calT_X \to \calT_X, \quad \partial \mapsto \partial^p$.
  \end{itemize}
  In characteristic $p$, the $p$-th power map of a derivation is still a derivation, which is a special phenomenon in positive characteristic.

  \begin{definition}
    Assume $\characteristic \kk = p > 0$.
    A \emph{foliation} on $X$ is a subbundle $\calF \subset \calT_X$ such that
    \begin{itemize}
      \item $\calF$ is involutive: $[\calF, \calF] \subset \calF$;
      \item $\calF$ is \emph{$p$-closed}: $\calF^{[p]} \subset \calF$.
    \end{itemize}
  \end{definition}

  \begin{example}
    Let $X = \bbA^n_\kk$ with $\characteristic \kk = p > 0$ and coordinates $x_1, \dots, x_n$.
    The foliation $\calF$ is generated by $\partial_{x_1}, \dots, \partial_{x_r}$, where $r < n$.
    Consider the differential equation
    \[ \partial_{x_1} f = 0, \dots, \partial_{x_r} f = 0 \]
    for $f \in \kk[[x_1, \dots, x_n]]$.
    The solution is 
    \[ \kk[x_1^p, \dots, x_r^p, x_{r+1}, \dots, x_n]. \]
    We get a morphism 
    \[ \Spec \kk[x_1^p, \dots, x_r^p, x_{r+1}, \dots, x_n] \to \bbA^n_\kk \]
    instead of the expected $\Spec \kk[x_{r+1}, \dots, x_n] \to \bbA^n_\kk$.
    The result morphism is a finite, flat, purely inseparable morphism of degree $p^r$.
  \end{example}

  \begin{theorem}[{\cite{Miyaoka-Peternell, Prop 1.9}}]
    Let $\kk$ be a field of characteristic $p > 0$, $X/\kk$ a smooth variety, and $\calF$ a foliation on $X$.
    Then there exists a canonical bijection between the following two sets:
    \begin{itemize}
      \item the set of foliations $\calF$ on $X$;
      \item smooth varieties $Y/\kk$ between $X$ and $X^{(p)}$ (i.e. $X \xrightarrow{f} Y \xrightarrow{g} X^{(p)}$).
    \end{itemize}
    Given a foliation $\calF$, the corresponding $Y$ is $\Spec \calO_X^\calF$, where for any affine open $U \subset X$,
    \[ \calO_X^\calF(U) = \{ f \in \calO_X(U) : \partial f = 0, \forall \partial \in \calF(U) \}. \]
    Conversely, given $Y$ between $X$ and $X^{(p)}$, the corresponding foliation $\calF$ is given by $T_{X/Y} = \Ker(\calT_X \to f^* \calT_Y)$.
  \end{theorem}

\subsection{Foliation over number fields}

  Let $\kk$ be a number field, $X/\kk$ a smooth variety, and $\calF$ a foliation on $X$.
  We ask the algebraicity of formal leaves, or somehow the algebraicity of power series.

  \begin{theorem}[{\cite{Borel}}]
    Let $f(t) \in \bbZ[[t]]$ is a power series with integer coefficients.
    If the radius of convergence of $f(t)$ over $\bbC$, $R_\infty(f)$, strictly larger than $1$,
    then $f(t)$ is a polynomial function. 
  \end{theorem}

  The condition $f \in \bbZ[[t]]$ is equivalent to that the $p$-adic radius of convergence, $R_p(f)$, is at least $1$ for all $p$.


  \begin{conjecture}\label{conj:algebraicity_of_leaves_by_reduction}
    Let $\kk$ be a number field, $X/\kk$ a smooth variety, $\calF$ a foliation on $X$.
    Suppose that $\calF$ is $p$-closed for almost all $\frakp \subset \calO_\kk$.
    Then there exists a Zariski open subset $U \subset X$ such that every formal leaf through a point in $U(\overline{\kk})$ is algebraic.
  \end{conjecture}

  Essentially, the fact that leaves are algebraic implies that the foliation comes from a fibration.

  We have the following partial result by Bost.

  \begin{theorem}[{\cite{Bost}}]\label{thm:Bost_main_theorem}
    In the \cref{conj:algebraicity_of_leaves_by_reduction}, if for some point $P \in X(\kk)$, the holomorphic leaf $L_P$ through $P$ for some embedding $\sigma: \kk \hookrightarrow \bbC$ satisfies the \emph{Liouville condition}, then the formal leaf through $P$ is algebraic (as subvariety of $X$).
  \end{theorem}

  \begin{definition}
    Let $M$ be a complex manifold.
    The \emph{Liouville condition} says that all bounded p.s.h. functions on $M$ are constant.
  \end{definition}

  \begin{remark}
    If $M$ is a complex algebraic variety, then $M$ satisfies the Liouville condition.
  \end{remark}

  \begin{example}
    Let $X$ be a abelian variety, $\calF \subset \calT_X$ be a foliation which is translation invariant (i.e. $\calF$ is generated by global sections of $\calT_X$).
    Then all holomorphic leaves are image of $\bbC^r$.
    Hence all holomorphic leaves satisfy the Liouville condition.
  \end{example}

  \begin{corollary}\label{cor:isogeny_of_elliptic_curves_by_foliation}
    Let $E_i$ be elliptic curves over $\bbQ$ for $i = 1, 2$.
    TFAE:
    \begin{enumerate}
      \item $E_1$ and $E_2$ are isogenous;
      \item $a_p(E_1) = a_p(E_2)$ for almost all $p$.
    \end{enumerate}
  \end{corollary}
  Recall $a_p(E) = p + 1 - \#E(\bbF_p) = \tr(\Frob_p|T_\ell(E))$ for an elliptic curve $E$ over $\bbQ$.
  If $E$ is an elliptic curve over $\bbF_p$, $D \in \Gamma(E, T_E)$ is a nonzero global section, then $D^p = a_p(E) D$.
  \begin{proof}
    Note that (a) $\Rightarrow$ (b) is trivial.
    For (b) $\Rightarrow$ (a), take $X = E_1 \times E_2$ and $D = p^*_1 D_1 + p^*_2 D_2$, where $D_i \in \Gamma(E_i, T_{E_i})$ is a nonzero global section.
    Then $D$ generates a foliation $\calF$ on $X$.
    For almost all $p$, 
    \[ \calD^p = p^*_1 (D_1)^p + p^*_2 (D_2)^p = a_p(E_1) p^*_1 D_1 + a_p(E_2) p^*_2 D_2 = a_p (p^*_1 D_1 + p^*_2 D_2) \mod p. \]
    Hence $\calF$ is $p$-closed for almost all $p$.
    By \cref{thm:Bost_main_theorem}, we are done.
  \end{proof}

  \begin{remark}
    The \cref{cor:isogeny_of_elliptic_curves_by_foliation} is originally proved by Faltings.
    It is a special case of the following theorem.
  \end{remark}

  \begin{theorem}[Faltings, Tate conjecture]
    Let $A_1$ and $A_2$ be abelian varieties over a number field $\kk$.
    The homomorphism
    \[ \Hom(A_1, A_2) \otimes_\bbZ \bbQ_\ell \to \Hom_{\Gal(\overline{\kk}/\kk)}(T_\ell(A_1), T_\ell(A_2)) \]
    is an isomorphism.
  \end{theorem}

  \begin{remark}
    In the \cref{conj:algebraicity_of_leaves_by_reduction}, we get a formal leaf $\hat{L}_P$ through $P$.
    The $p$-closed condition implies the ``radius of convergence'' $R_p$ of $\hat{L}_P$ has a lower bound only depending on $p$.
    Moreover, we have 
    \[ \prod_{p \text{ good }} R_p > 0. \]
    \textbf{Roughly}, we can view the Liouville condition implies that the radius of convergence $R_\infty$ of $\hat{L}_P$ over $\bbC$ is infinite.
  \end{remark}

  \paragraph{Grothendieck-Katz $p$-curvature conjecture}

  We consider more conjectures.

  \begin{conjecture}[Grothendieck]\label{conj:ODE_reduction}
    Let $\kk$ be a number field.
    Consider an ODE for $X = (x_1(t), \dots, x_n(t))^T \in \overline{\kk(t)}^n$:
    \[ \frac{d X}{d t} = - A(t) X \]
    where $A(t) \in M_n(\kk(t))$ is a matrix with rational function entries. 
    TFAE:
    \begin{enumerate}
      \item the ODE has $n$ $\overline{\kk(t)}$-linearly independent solutions in $\overline{\kk(t)}^n$;
      \item for almost all $\frakp \subset \calO_\kk$, the reduction of the ODE modulo $\frakp$ has $n$ $\overline{\bbF_p(t)}$-linearly independent solutions over $\overline{\bbF_p(t)}^n$.
    \end{enumerate}
  \end{conjecture}

  Let $\kk$ be a  field, $X/\kk$ a smooth variety, $\calE$ a vector bundle on $X$ with a connection $\nabla: \calE \to \calE \otimes \Omega^1_X$.
  Recall that $\nabla$ is a connection means that $\nabla$ satisfies the Leibniz rule: $\nabla(f s) = f \nabla(s) + s \otimes df$ for any $f \in \calO_X(U)$ and $s \in \calE(U)$.

  \begin{definition}
    \begin{enumerate}
      \item $\calE^\nabla$ is the subsheaf of $\calE$ consisting of horizontal sections, i.e. $\calE^\nabla(U) = \{ s \in \calE(U) : \nabla(s) = 0 \}$ for any affine open $U \subset X$.
      \item we say $\nabla$ is \emph{trivial} if $\calE = \calE^\nabla \otimes_\kk \calO_X$ (generated by horizontal sections).
       we say $\nabla$ is \emph{isotrivial} if there exists a finite étale cover $f: Y \to X$ such that $f^* \nabla$ is trivial.
      \item say $(\calE, \nabla)$ is \emph{integrable} if the \emph{curvature} $\nabla^2: \calE \to \calE \otimes \Omega^2_X$ is zero, i.e. $\nabla_{[u,v]} - [\nabla_u, \nabla_v] = 0 \in \End_\kk(\calE)$ for any vector fields $u, v$ on $X$.
        Here we view $\nabla$ as a map $\calT_X \to \End_\kk(\calE)$ by $\nabla_u(s) = \nabla(s)(u)$ for any $u \in \calT_X(U)$ and $s \in \calE(U)$.
      \item if $\characteristic \kk = p > 0$, we define the $p$-curvature 
        \[ \psi_p: \calT_X \to \End_\kk(\calE), \quad u \mapsto \nabla_u^p - \nabla_{u^{[p]}}. \]
    \end{enumerate}
  \end{definition}

  \begin{conjecture}[Katz]\label{conj:p_curvature}
    Let $\kk$ be a number field, $X/\kk$ a smooth variety, and $(\calE, \nabla)$ an integrable connection on $X$.
    TFAE:
    \begin{enumerate}
      \item $(\calE, \nabla)$ is isotrivial;
      \item for almost all $\frakp \subset \calO_\kk$, the reduction of $(\calE, \nabla)$ modulo $\frakp$ has zero $p$-curvature.
    \end{enumerate}
  \end{conjecture}

  \cref{conj:ODE_reduction,conj:p_curvature} are equivalent.

  \cref{conj:algebraicity_of_leaves_by_reduction} can imply \cref{conj:p_curvature}.

  Take the geometric vector bundle $W := \bbV(\calE) = \Spec \Sym^\bullet \calE$ over $X$ via $\pi: W \to X$.
  For any $v \in T_X(U)$, we have $\nabla_v: \calE(U) \to \calE(U)$ acting on $\Sym^\bullet \calE(U)$ satisfying the Leibniz rule.
  Hence $\nabla_v$ is a derivation on $W$ since $\calO_W(U) = \Sym^\bullet \calE(U)$.
  This gives a morphism $\pi^*\calT_X \to \calT_W$.
  Denote by $\calF$ the image of $\pi^* \calT_X$ in $\calT_W$.
  Then $\calF$ is a foliation on $W$.
  